\SourceLeaf{217}
Werth von \(\psi\) verschwinden, und in allen Punkten des Raumes, welche der Fläche nicht angehören, wird er nicht verschwinden. Bei einem Durchgange durch die Fläche wird also, wenigstens allgemein zu reden, der Werth von \(\psi\) aus dem Positiven ins Negative, bei dem entgegengesetzten aus dem Negativen ins Positive übergehen, oder auf der einen Seite der Fläche wird der Werth von \(\psi\) positiv, auf der andern negativ sein: die erstere wollen wir als die obere, die andere als die untere betrachten. Ganz eben so soll es bei der zweiten Fläche gehalten werden, indem ihre Natur durch die Gleichung \(\Psi=0\) bestimmt wird, wo \(\Psi\) eine gegebne einförmige Function der Coordinaten \(X,Y,Z\) ist. Es gebe ferner die Differentiation
\[
\begin{aligned}
\mathrm{d}\psi&=e\,\mathrm{d}x+g\,\mathrm{d}y+h\,\mathrm{d}z\\
\mathrm{d}\Psi&=E\,\mathrm{d}X+G\,\mathrm{d}Y+H\,\mathrm{d}Z
\end{aligned}
\]
\noindent wo \(e,g,h\) Functionen von \(x,y,z\) und \(E,G,H\) Functionen von \(X,Y,Z\) sein werden.

Da die Betrachtungen, durch welche wir zu dem vorgesetzten Ziele gelangen müssen, obwohl an sich nicht schwierig, doch etwas ungewöhnlicher Art sind, so wollen wir uns bemühen, ihnen die grösste Klarheit zu geben. Wir wollen zwischen den beiden einander entsprechenden Darstellungen auf dem Flächen, deren Gleichungen \(\psi=0\) und \(\Psi=0\) sind, sechs Zwischen-Darstellungen in der Ebene annehmen, so dass acht verschiedene Darstellungen in Betracht kommen, nemlich
\begin{center}
\begingroup\fontsize{8}{10}\selectfont\setlength{\tabcolsep}{2pt}
\begin{tabular}{@{}lll@{}}
\multicolumn{3}{r}{\parbox{1.55in}{\raggedright indem als correspondirend betrachtet werden die Punkte, deren Coordinaten resp. \(=\)}}\\[2pt]
1\textsuperscript{s}&das Urbild in der Fläche, deren Gleichung \(\psi=0\) . .&\(x,\ y,\ z\)\\
2\textsuperscript{s}&Darstellung in der Ebene . . . . . . . . . . . . . . . . . . . . .&\(x,\ y,\ 0\)\\
3\textsuperscript{s}&\quad „\quad „\quad „\quad „\quad . . . . . . . . . . . . . . . . . . . . .&\(t,\ u,\ 0\)\\
4\textsuperscript{s}&\quad „\quad „\quad „\quad „\quad . . . . . . . . . . . . . . . . . . . . .&\(p,\ q,\ 0\)\\
5\textsuperscript{s}&\quad „\quad „\quad „\quad „\quad . . . . . . . . . . . . . . . . . . . . .&\(P,\ Q,\ 0\)\\
6\textsuperscript{s}&\quad „\quad „\quad „\quad „\quad . . . . . . . . . . . . . . . . . . . . .&\(T,\ U,\ 0\)\\
7\textsuperscript{s}&\quad „\quad „\quad „\quad „\quad . . . . . . . . . . . . . . . . . . . . .&\(X,\ Y,\ 0\)\\
8\textsuperscript{s}&Abbildung in der Fläche, deren Gleichung \(\Psi=0\) . .&\(X,\ Y,\ Z\)
\end{tabular}\endgroup
\end{center}
\par\indent Wir wollen nun diese verschiedenen Darstellungen unter einander lediglich in Beziehung auf die gegenseitige \textit{Lage} der unendlich kleinen Linearelemente ver%
